\section{Experimental Setups}
\subsection{Dataset}
We construct our synthetic training set using the LLaVA-OneVision~\cite{li2024llavaonevisioneasyvisualtask}, which covers diverse multimodal tasks including reasoning, math and coding, and captioning. To focus on multimodal evaluation, we restrict ourselves to the single-image subsets.
To avoid dominance by any single sub-dataset, we impose a cap of 10k on the number of examples drawn from each sub-dataset, resulting in 100k prompts. This ensures balanced coverage across categories while maintaining manageable data size. 


\begin{table*}[t]
\centering
\small
\setlength{\tabcolsep}{6pt}
\begin{tabular}{lcccccccccccc}
\toprule
 & \multicolumn{4}{c}{\textit{VLRB}} & \multicolumn{6}{c}{\textit{MMRB}} \\
\cmidrule(lr){2-5} \cmidrule(lr){6-11}
\textbf{Model} & \cellcolor{gray!20}\textbf{Ave.} & \textbf{Gen.} & \textbf{Hallu.} & \textbf{Reason.} & \cellcolor{gray!20}\textbf{Ave.} & \textbf{Gen.} & \textbf{Know.} & \textbf{Reason.} & \textbf{Safe.} & \textbf{VQA} \\
\midrule
\multicolumn{11}{l}{\textit{Larger Models}} \\
Claude-3.5 & \cellcolor{gray!20}0.536 & 0.434 & 0.550 & 0.623 & \cellcolor{gray!20}0.721 & 0.652 & 0.739 & 0.669 & 0.687 & 0.856 \\
GPT-4o & \cellcolor{gray!20}0.624 & 0.491 & 0.676 & 0.705 & \cellcolor{gray!20}0.713 & 0.658 & 0.720 & 0.649 & 0.668 & 0.872 \\
Llama-90B & \cellcolor{gray!20}0.539 & 0.426 & 0.573 & 0.617 & \cellcolor{gray!20}0.618 & 0.642 & 0.612 & 0.547 & 0.519 & 0.771 \\
\midrule
\multicolumn{11}{l}{\textit{Our Method (Based on Llama-3.2-11B-Vision-Instruct)}} \\
Base & \cellcolor{gray!20}0.383 & 0.297 & 0.304 & 0.549 & \cellcolor{gray!20}0.499 & 0.590 & 0.506 & 0.517 & 0.317 & 0.565 \\
Iteration 1 & \cellcolor{gray!20}0.452 & 0.398 & 0.362 & 0.596 & \cellcolor{gray!20}0.519 & \textbf{0.625} & 0.532 & 0.500 & 0.289 & 0.651 \\
Iteration 2 & \cellcolor{gray!20}\underline{0.521} & \underline{0.453} & \textbf{0.529} & 0.580 & \cellcolor{gray!20}0.521 & 0.591 & 0.506 & 0.513 & 0.304 & \underline{0.693} \\
Iteration 3 & \cellcolor{gray!20}0.488 & 0.425 & 0.426 & \textbf{0.612} & \cellcolor{gray!20}\underline{0.538} & \underline{0.599} & \underline{0.540} & \textbf{0.543} & 0.307 & \textbf{0.701} \\
Iteration 4 & \cellcolor{gray!20}\textbf{0.538} & \textbf{0.503} & \underline{0.514} & \underline{0.596} & \cellcolor{gray!20}\textbf{0.539} & 0.591 & \textbf{0.556} & \underline{0.531} & \textbf{0.329} & 0.689 \\
\bottomrule
\end{tabular}
\caption{Performance comparison on VLRB and MMRB benchmarks. Our iteratively trained judge model based on Llama-3.2-11B achieves competitive performance with significantly larger models across multiple evaluation dimensions.}
\label{tab:vlrb_mmrb_results}
\end{table*}

\subsection{Evaluation}

We evaluate our VLM judge and baselines on two comprehensive benchmarks: Multimodal RewardBench~\cite{mmrb} and VL-RewardBench~\cite{vlrb}. Both benchmarks measure judge performance across multiple dimensions including general instructions, hallucinations, correctness, reasoning, and VQA. Multimodal RewardBench comprises 5,211 preference pairs covering both long-form and short-answer formats, with labels derived from ground truth correctness or high-quality human expert annotations. VL-RewardBench spans general multimodal instructions, hallucination detection, and complex reasoning tasks, combining human-verified and AI-annotated preference labels to ensure comprehensive coverage of real-world scenarios.

\subsection{Implementation details}

We fine-tune the Llama-3.2-11B-Vision-Instruct model for 5 epochs using a learning rate of 1e-5 and a batch size of 2 per GPU across 8 GPUs (effective batch size of 16). We employ FSDP for distributed training and enable fast kernels for computational efficiency. We use cross-entropy loss computed over the generated reasoning traces from each iteration cycle. 
We iterate the training process until overall performance plateaus, operationally 
defined as less than 1\% relative improvement over three consecutive iterations on either of the benchmarks. 
In our experiments, we observe diminishing returns after iteration 4, though individual dimensions may peak at different iterations due to varying learning 
dynamics across task types. More details on the hyperparameters in data generation, sampling, and training can be found in Appendix~\ref{app:hyperparameters}.




\section{Results}

\subsection{Iterative Training Improvements} Table~\ref{tab:vlrb_mmrb_results} demonstrates the effectiveness of iterative training on our Llama3-11B based judge model. On VLRB, our model improves from 0.383 (base) to 0.538 (iteration 4), representing a 40.5\% relative gain. On MMRB, we observe improvement from 0.499 (base) to 0.539 (iteration 4), a 7.5\% relative gain. 
The progression across iterations demonstrates relatively consistent improvement, though the improvement trajectory can vary across different dimensions.

Compared with larger models, despite using only 11B parameters, our model achieves competitive performance on several dimensions: on general instruction following, our model (0.503) substantially outperforms the 90B parameter Llama-3.2-90B (0.426), Claude-3.5-Sonnet (0.434) and GPT-4o (0.491). Across VLRB hallucination and MMRB VQA, our model (0.514, 0.689) approaches larger models despite the base model's significantly weaker starting performance (0.304, 0.565).

\subsection{Performance Gains Across Dimensions}

Our iterative training exhibits dimension-specific learning dynamics that reveal the differential effectiveness of our synthetic data generation strategies. On VLRB, general instruction following demonstrates the largest relative gain of 69\% (0.297 $\rightarrow$ 0.503), indicating that our preference pair construction successfully captures the supervision signal necessary for open-ended multimodal instruction assessment. Hallucination detection on VLRB achieves 40.9\% relative improvement (0.304 $\rightarrow$ 0.514), validating our detail alteration methodology for training factual consistency discrimination. Reasoning evaluation shows 8.6\% relative improvement (0.549 $\rightarrow$ 0.596 at iteration 4), though exhibits non-monotonic optimization behavior with peak performance at iteration 3 (0.612), suggesting potential signs of overfitting in later iterations.

On MMRB, VQA tasks yield substantial gains (18.0\% relative improvement, 0.565 $\rightarrow$ 0.689), confirming effective transfer from our majority voting-based preference construction for closed-ended answers. In contrast, several dimensions show relatively marginal improvement: reasoning capability increases from 0.517 to 0.531 (2.7\% relative gain), general evaluation (averaging correctness and preference) stayed flat from 0.590 $\rightarrow$ 0.591, and safety evaluation demonstrates limited responsiveness (0.317 $\rightarrow$ 0.329). These marginal gains indicate potential distributional mismatch between our synthetic data and the supervision required for these specific capabilities.

